\section{Architecture}
\label{sec:arch}

\subsection{AI Chain is the Lux A-Chain}

Hanzo AI Chain is the public identity of the Lux A-Chain when it operates the
open AI protocol. There is not a Hanzo proof chain beside A-Chain. A single
canonical state machine owns job admission, proof policy, challenge state, work
nullifiers, reward authorization, and receipt roots.

\begin{figure}[h]
\centering
\resizebox{\linewidth}{!}{%
\begin{tikzpicture}[
    node distance=0.8cm and 0.85cm,
    box/.style={rectangle, draw, minimum width=2.3cm, minimum height=0.7cm,
                font=\small, rounded corners=1pt},
    arrow/.style={->, thick}
]
\node[box] (job) {Demand-backed AI job};
\node[box, right=of job] (exec) {CPU/GPU execution};
\node[box, right=of exec] (a) {A-Chain PoAI};
\node[box, right=of a] (z) {Z-Chain P3Q};
\node[box, right=of z] (receipt) {Settled receipt};
\node[box, below=of z] (lux) {Lux Quasar + PQ policy};
\node[box, below=of receipt] (dest) {Bound destination};
\draw[arrow] (job) -- (exec);
\draw[arrow] (exec) -- node[above,font=\scriptsize]{commit + proof} (a);
\draw[arrow] (a) -- (z);
\draw[arrow] (z) -- (receipt);
\draw[arrow] (lux) -- (z);
\draw[arrow] (receipt) -- (dest);
\end{tikzpicture}
}
\caption{PoAI validates useful work on A-Chain; P3Q settles it on Z-Chain; Lux
supplies PQ-rooted finality; destinations consume receipts rather than raw work.}
\label{fig:arch}
\end{figure}

\subsection{PoAI useful-work state machine}

Each job progresses through

\begin{equation}
\begin{aligned}
\textsc{Requested}&\rightarrow\textsc{Committed}
\rightarrow\textsc{Verifying/Challengeable}\\
&\rightarrow\textsc{Finalized}\;|\;\textsc{Rejected}\;|\;\textsc{Expired}.
\end{aligned}
\end{equation}

A-Chain fixes the execution and proof policy before commitment. It derives any
random challenge afterward, resolves the policy to final or rejected, checks
metering and reward allowance, consumes the global nullifier, and creates the
destination-bound receipt atomically. An optimistic result cannot become
issuance-final while its challenge window remains open.

\subsection{Z-Chain settlement and Lux base-layer finality}

Z-Chain is the P3Q proof-compression and settlement rollup for finalized A-Chain
receipt transitions. Its intended proof system is a Plonky3-derived,
pairing-free STARK/FRI stack. Z-Chain does not admit jobs, determine usefulness,
own the mining nullifier, or authorize subsidy; its proof must bind exactly to
the finalized A-to-Z receipt transition.

Lux Quasar orders and finalizes A-Chain and Z-Chain blocks. Exported roots must be
authenticated under the validator and post-quantum policy that is actually
activated. The security statement therefore names its assumptions: validator
set, threshold, key lifecycle, downgrade prevention, and relay verification.
Selecting ML-DSA in a registry does not by itself produce a threshold finality
certificate.

Pulsar and Corona explore threshold lattice signatures for consensus and
custody. P3Q is the intended Z-Chain settlement substrate, but the exact PoAI
AIR, recursion policy, state-transition binding, and verifier remain subject to
protocol, audit, integration, and activation gates. This paper does not treat a
repository or precompile address as deployment evidence.

\subsection{Deterministic execution substrate}

The public proof-bearing path admits a bounded deterministic graph whose
operations define:

\begin{itemize}[leftmargin=1.2em]
  \item integer weights and activations, normally \texttt{int8};
  \item wider accumulators with explicit overflow or saturation behavior;
  \item canonical tensor layout, serialization, reduction, and rounding;
  \item deterministic routing, normalization, sampling, and tie-breaking;
  \item content-addressed model, input, output, runtime, and graph roots.
\end{itemize}

Execution may occur on a CPU, GPU, NPU, or other backend that produces the same
committed bytes. Accelerators improve throughput; they are not a consensus
requirement for every validator.

\subsection{Evidence backends}

The job selects one activated evidence profile:

\begin{table}[h]
\centering
\small
\begin{tabular}{p{0.23\linewidth}p{0.67\linewidth}}
\toprule
\textbf{Profile} & \textbf{Role} \\
\midrule
Exact replay & Canonical CPU or verifier-set recomputation for small jobs. \\
Optimistic integer & Transcript commitment, data availability, watcher
re-execution, and objective challenge openings. \\
Succinct P3Q & Plonky3-derived Z-Chain settlement only after the exact PoAI AIR,
A-to-Z binding, recursion policy, and verifier are audited and activated. \\
Confidential compute & Optional vendor attestation for private prompts, weights,
or outputs, combined with explicit TCB and revocation policy. \\
Subjective quorum & Attributable judgment that may earn demand fees but cannot
alone authorize mining subsidy. \\
\bottomrule
\end{tabular}
\caption{Evidence profiles are selected by job policy; none creates another
mining authority outside A-Chain.}
\end{table}

\subsection{Canonical protocol objects}

The minimum consensus objects are \texttt{Mining\allowbreak Job},
\texttt{Execution\allowbreak Commitment}, \texttt{Challenge},
\texttt{Proof\allowbreak Finality}, \texttt{Work\allowbreak Nullifier}, and
\texttt{Mining\allowbreak Receipt}. Public APIs are submission
and query surfaces for these objects. An in-memory task server, GPU scheduler,
or destination smart contract is not the canonical mining state machine.

\subsection{Implementation status}

The exact-integer arithmetic core, Go/Rust vectors, Solidity witness fixtures,
and A-Chain receipt components exist as implementation evidence. The production
system still requires a complete Zen model transcript, data-availability path,
finalized-not-pending proof gate, verifiable metering, token-policy upgrade, PQ
root export, Z-Chain settlement, and exactly-once destination consumption. Mining remains disabled
until the launch gates in \S\ref{sec:deploy} pass.
